\documentclass[a4paper]{arrowhead}

\usepackage[yyyymmdd]{datetime}
\usepackage{etoolbox}
\usepackage[utf8]{inputenc}
\usepackage[T1]{fontenc}
\usepackage{multirow}
\usepackage{longtable}
\usepackage{tikz}
\usetikzlibrary{positioning,shapes.multipart,arrows.meta}

\renewcommand{\dateseparator}{-}
\setlength{\emergencystretch}{3em}
\lstset{language=HTTP,breakatwhitespace=false,columns=fullflexible}

%% Special references
\newcommand{\mref}[2][]{{\textcolor{ArrowheadPurple}{\hyperref[sec:model:#2]{\ifstrempty{#1}{#2}{#1}}}}}
\newcommand{\oref}[1]{{\textcolor{ArrowheadBlue}{\hyperref[sec:operations:#1]{#1}}}}
\newcommand{\pdef}[1]{{\textcolor{ArrowheadGrey}{#1\label{sec:model:primitives:#1}\label{sec:model:primitives:#1s}\label{sec:model:primitives:#1es}}}}
\newcommand{\pref}[1]{{\textcolor{ArrowheadGrey}{\hyperref[sec:model:primitives:#1]{#1}}}}

% as template_sd.tex; the input and output arguments are passed already formatted (\mref or \pref)
\newrobustcmd\osubsection[3]{
  \addtocounter{subsection}{1}
  \addcontentsline{toc}{subsection}{\protect\numberline{\thesubsection}operation \textcolor{ArrowheadBlue}{#1}}
  \subsection*{
    \thesubsection\quad
    operation
    \textcolor{ArrowheadBlue}{#1}
    (\notblank{#2}{#2}{})
    \notblank{#3}{: #3}{}
  }
  \label{sec:operations:#1}
}
\newrobustcmd\msubsection[2]{
  \addtocounter{subsection}{1}
  \addcontentsline{toc}{subsection}{\protect\numberline{\thesubsection}#1 \textcolor{ArrowheadPurple}{#2}}
  \subsection*{\thesubsection\quad#1 \textcolor{ArrowheadPurple}{#2}}
  \label{sec:model:#2} \label{sec:model:#2s} \label{sec:model:#2es}
}
%%

\usepackage{xurl}
\begin{document}
\sloppy

%% Arrowhead Document Properties
\ArrowheadTitle{deviceQualityDataManagement}
\ArrowheadServiceID{deviceQualityDataManagement}
\ArrowheadType{Service Description}
\ArrowheadTypeShort{SD}
\ArrowheadVersion{5.3.0}
\ArrowheadDate{2026-11-09}
\ArrowheadAuthor{Jerker Delsing}
\ArrowheadStatus{REVIEW DRAFT}
\ArrowheadContact{jerker.delsing@ltu.se}
\ArrowheadFooter{\href{www.arrowhead.eu}{www.arrowhead.eu}}
\ArrowheadSetup
%%

%% Front Page
\begin{center}
  \vspace*{1cm}
  \huge{\arrowtitle}

  \vspace*{0.2cm}
  \LARGE{\arrowtype}
  \vspace*{1cm}

  \Large{Service ID: \textit{"\arrowid"}}

  \vspace*{\fill}

  % Front Page Image

  \vspace*{1cm}
  \vspace*{\fill}

  % Front Page Abstract
  \begin{abstract}
    This document describes the \texttt{device\allowbreak{}Quality\allowbreak{}Data\allowbreak{}Management} service, version 1.0.0, produced by the DeviceQoSEvaluator support system of Arrowhead (documentation edition 5.3.0).
  \end{abstract}

  \vspace*{1cm}
\end{center}

\newpage
%%

%% Table of Contents
\tableofcontents
\newpage
%%

\section{Documentation edition and implementation scope}
This is the 5.3.0 documentation edition, dated 2026-11-09. It carries forward the system and service contracts from the retained 5.2.2 documentation and identifies the current primary implementation sources. The former \path|documentation/v5.2.2| folder is retained unchanged.

Document version 5.3.0 is distinct from executable, dependency, registered system/service and wire-contract versions. Values still defined by the source (including 5.2.2 artefact versions, 1.0.0 service versions and \path|SysMonitor-AMQP-IDD-5.2.2| payload identifiers) are preserved. This documentation update does not change runtime behaviour or establish a new implementation release. Historical verification statements describe the tests recorded in the inherited baseline; this revision validates document compilation and references.

\section{Overview}
\label{sec:overview}

This document describes the \texttt{device\allowbreak{}Quality\allowbreak{}Data\allowbreak{}Management} service, which provides operator access to the stored device statistics and rebuilding of the device-system measurement map.
It is produced by the DeviceQoSEvaluator system \cite{dqe-sysd} in service version 1.0.0.

The rest of this document is organized as follows.
In Section \ref{sec:operations}, we describe the abstract message operations provided by the service.
In Section \ref{sec:model}, we end the document by presenting the data types used by the mentioned operations.
The concrete interfaces are specified in the HTTP/TLS/JSON IDD \cite{dqe-deviceQualityDataManagement-http-idd} and the MQTT/TLS/JSON IDD \cite{dqe-deviceQualityDataManagement-mqtt-idd}.

\subsection{Significant Prior Art}

New in Arrowhead 5.

\subsection{How This Service Is Meant to Be Used}

Operators query aggregated statistics per metric group and period, and call \texttt{reload} after devices or systems changed in the ServiceRegistry.

\subsection{Important Delimitations and Dependencies}
\label{sec:delimitations}

Access follows the management policy of the Device\allowbreak{}Qo\allowbreak{}SEvaluator.

\newpage

\section{Service Interface}
\label{sec:operations}

This section describes the interfaces to the \texttt{device\allowbreak{}Quality\allowbreak{}Data\allowbreak{}Management} service.
In particular, each subsection names an abstract operation, an input type and an output type, in that order.
The input type is named inside parentheses, while the output type is preceded by a colon.
Input and output types are only denoted when accepted or returned, respectively, by the interface in question.

All abstract data types named in this section are defined in Section \ref{sec:model}.

\begin{figure}[ht!]
  \centering
  \begin{tikzpicture}[font=\small]
    \node[draw=ArrowheadBlue, very thick, rectangle split, rectangle split parts=2, align=center,
          rectangle split part fill={ArrowheadBlue!12,white}, text width=9cm, inner sep=6pt] (b)
      {\guillemotleft{}Service\guillemotright{}\\\textbf{deviceQualityDataManagement}
        \nodepart{second} \texttt{query}\\ \texttt{reload}};
  \end{tikzpicture}
  \caption{SysML-style block of the \texttt{device\allowbreak{}Quality\allowbreak{}Data\allowbreak{}Management} service and its operations.}
  \label{fig:SD}
\end{figure}

The following interface operations are available.

\osubsection{query}{\mref[Qo\allowbreak{}SDevice\allowbreak{}Stat\allowbreak{}Query\allowbreak{}Request]{QoSDeviceStatQueryRequest}}{\mref[Qo\allowbreak{}SDevice\allowbreak{}Stat\allowbreak{}Query\allowbreak{}Response]{QoSDeviceStatQueryResponse}}

Returns a page of aggregated statistic records.

\osubsection{reload}{}{\pref{String}}

Rebuilds the device-system measurement map from the ServiceRegistry.

\newpage

\section{Information Model}
\label{sec:model}

Here, all data objects that can be part of the \texttt{device\allowbreak{}Quality\allowbreak{}Data\allowbreak{}Management} service are listed in alphabetic order.
Note that each subsection, which describes one type of object, begins with the \textit{struct} keyword, which is used to denote a collection of named fields, each with its own data type.
As a complement to the explicitly defined types in this section, there is also a list of implicit primitive types in Section \ref{sec:model:primitives}, which are used to represent things like hashes and identifiers.

\msubsection{struct}{Page}

Paging and sorting; all fields optional.


\noindent\begin{tabularx}{\textwidth}{| >{\raggedright\arraybackslash}p{3.9cm} | >{\raggedright\arraybackslash\small}p{5.3cm} | X |} \hline
\rowcolor{gray!33} Field & Type & Description \\ \hline
page & \pref{Number} & Zero-based page number. \\ \hline
size & \pref{Number} & Number of records per page (at most \texttt{max.\allowbreak{}page.\allowbreak{}size}, default 1000). \\ \hline
direction & \pref{String} & Sort direction, \texttt{ASC} or \texttt{DESC}. \\ \hline
sort\allowbreak{}Field & \pref{String} & Field to sort by. \\ \hline
\end{tabularx}



\msubsection{struct}{QoSDeviceStatQueryRequest}

Paged filter for device statistics.


\noindent\begin{tabularx}{\textwidth}{| >{\raggedright\arraybackslash}p{3.9cm} | >{\raggedright\arraybackslash\small}p{5.3cm} | X |} \hline
\rowcolor{gray!33} Field & Type & Description \\ \hline
pagination & \mref{Page} & Paging and sorting of the result. \\ \hline
metric\allowbreak{}Group & \pref{String} & Metric group: \texttt{RTT}, \texttt{CPU\_\allowbreak{}TOTAL\_\allowbreak{}LOAD}, \texttt{MEMORY\_\allowbreak{}USED}, \texttt{NETWORK\_\allowbreak{}EGRESS\_\allowbreak{}LOAD} or \texttt{NETWORK\_\allowbreak{}INGRESS\_\allowbreak{}LOAD}. \\ \hline
from & \pref{DateTime} & Lower bound of the entry time. \\ \hline
to & \pref{DateTime} & Upper bound of the entry time. \\ \hline
aggregation & \pref{List}$<$\allowbreak{}\pref{String}$>$ & Statistics to return (minimum, maximum, mean, median, current). \\ \hline
system\allowbreak{}Names & \pref{List}$<$\allowbreak{}\pref{Name}$>$ & Systems whose devices are returned. \\ \hline
\end{tabularx}



\msubsection{struct}{QoSDeviceStatQueryResponse}

Device statistic records.


\noindent\begin{tabularx}{\textwidth}{| >{\raggedright\arraybackslash}p{3.9cm} | >{\raggedright\arraybackslash\small}p{5.3cm} | X |} \hline
\rowcolor{gray!33} Field & Type & Description \\ \hline
records & \pref{List}$<$\allowbreak{}\mref{QoSDeviceStatRecord}$>$ & Statistic records of the page. \\ \hline
count & \pref{Number} & Total number of matching records (all pages). \\ \hline
\end{tabularx}



\msubsection{struct}{QoSDeviceStatRecord}

Aggregated measurements of one device in one period.


\noindent\begin{tabularx}{\textwidth}{| >{\raggedright\arraybackslash}p{3.9cm} | >{\raggedright\arraybackslash\small}p{5.3cm} | X |} \hline
\rowcolor{gray!33} Field & Type & Description \\ \hline
metric\allowbreak{}Group & \pref{String} & Metric group: \texttt{RTT}, \texttt{CPU\_\allowbreak{}TOTAL\_\allowbreak{}LOAD}, \texttt{MEMORY\_\allowbreak{}USED}, \texttt{NETWORK\_\allowbreak{}EGRESS\_\allowbreak{}LOAD} or \texttt{NETWORK\_\allowbreak{}INGRESS\_\allowbreak{}LOAD}. \\ \hline
id & \pref{Number} & Identifier of the record. \\ \hline
timestamp & \pref{String} & End of the aggregation period. \\ \hline
uuid & \pref{String} & Identifier of the measured device. \\ \hline
minimum & \pref{String} & Minimum value in the period. \\ \hline
maximum & \pref{String} & Maximum value in the period. \\ \hline
mean & \pref{String} & Mean value in the period. \\ \hline
median & \pref{String} & Median value in the period. \\ \hline
current & \pref{String} & Last measured value. \\ \hline
systems & \pref{List}$<$\allowbreak{}\pref{String}$>$ & Systems hosted by the device. \\ \hline
\end{tabularx}



\subsection{Primitives}
\label{sec:model:primitives}

Types and structures mentioned throughout this document that are assumed to be available to implementations of this service.
The concrete interpretations of each of these types and structures must be provided by any IDD document claiming to implement this service.

\noindent\begin{tabularx}{\textwidth}{| p{4.6cm} | X |} \hline
\rowcolor{gray!33} Type & Description \\ \hline
\pdef{DateTime} & Pinpoints a specific moment in time (UTC). \\ \hline
\pdef{List} & An \textit{array} of a known number of items, each having type A. \\ \hline
\pdef{Name} & A string identifier that is intended to be both human and machine-readable. \\ \hline
\pdef{Number} & An integer number. \\ \hline
\pdef{String} & A chain of characters. \\ \hline
\end{tabularx}

\newpage

\bibliographystyle{IEEEtran}
\bibliography{bibliography}

\newpage

\section{Revision History}
\subsection{Amendments}

\noindent\begin{tabularx}{\textwidth}{| p{1cm} | p{3cm} | p{2cm} | X | p{4cm} |} \hline
\rowcolor{gray!33} No. & Date & Version & Subject of Amendments & Author \\ \hline

1 & 2026-10-22 & 5.2.2 & Document created from the Arrowhead LaTeX templates for the 5.2.2 Java Spring implementation, including the local extensions of this workspace. & Jerker Delsing \\ \hline
2 & 2026-10-30 & 5.2.2 & Checked against the primary Java source tree (Eclipse-Arrowhead-java): services, operations, paths and data models unchanged. & Claude Opus 5.5 (draft) \\ \hline
3 & 2026-11-09 & 5.3.0 & New documentation edition; primary-source references and Go/Rust implementation scope; 5.2.2 retained. & Documentation maintenance \\ \hline
\end{tabularx}

\subsection{Quality Assurance}

\noindent\begin{tabularx}{\textwidth}{| p{1cm} | p{3cm} | p{2cm} | X |} \hline
\rowcolor{gray!33} No. & Date & Version & Approved by \\ \hline

1 & & \arrowversion  &  \\ \hline

\end{tabularx}

\end{document}
